\begin{enumerate}
    \item{\textbf{Calcula} la solución óptima de la relajación lineal de este problema:
    
\begin{equation}
\begin{split}
\mbox{max. } z = 5x_1 + 8x_2\\
s.a.:\\
x_1 + x_2 \leq  6\\
5x_1 + 9x_2 \leq 45\\
x_1,\,\,x_2 \in  \mathcal{Z}^+
\end{split}
\label{modeloGomory}
\end{equation}

 La solución óptima de la relación lineal de \eqref{modeloGomory} es:

\begin{center}
\begin{tabular}{c|c|cccc|}
 & $z$ & $x_1$ & $x_2$ &  $h_1$ & $h_2$ \\ 
 & -165/4 & 0 & 0  & -5/4 & -3/4\\ 
\hline
$x_1$ & 9/4  & 1 & 0 & 9/4 & -1/4\\ 
$x_2$ & 15/4  & 0 & 1 & -5/4 & 1/4\\ 
\hline
\end{tabular}
\end{center}

    \item{\textbf{}{Indica} la expresión del corte de Gomory asociado a la variable que tiene parte fraccional de mayor valor en la solución óptima de la relajación lineal obtenida.}

   $\frac{3}{4} h_1 + \frac{1}{4} h_2 \geq \frac{3}{4}$

    
    \item {Añade el corte propuesto y \textbf{obtén} la solución óptima del problema original, problema \eqref{modeloGomory}}.

    La solución óptima es $x^*_{PE}=(0,5)$ con $z^*_{PE}=40$
    \item ¿Cuál es la \textbf{expresión del corte introducido} en función de $x_1$ y $x_2$?
    
    $2x_1+3x_2 \leq 15$
    

\end{enumerate}